\documentclass[11pt]{article}
\usepackage[margin=1in]{geometry}
\usepackage{amsmath,amssymb,amsthm}
\usepackage{graphicx}
\usepackage[hidelinks]{hyperref}
\setlength{\emergencystretch}{2em}

\newtheorem{theorem}{Theorem}
\newtheorem{lemma}{Lemma}
\newtheorem{corollary}{Corollary}
\newtheorem{remark}{Remark}

\newcommand{\C}{\mathbb C}
\newcommand{\T}{\mathbb T}
\newcommand{\Bcal}{\mathcal B}
\newcommand{\Pcal}{\mathcal P}
\newcommand{\ext}{\operatorname{ext}}

\title{Pure Transition Pi-Bases Calibrate the Boundary Cases}
\author{Run \texttt{fa\_banach\_001}, lane 19}
\date{June 17, 2026}

\begin{document}
\maketitle

\section*{Source Problem}

The source paper is Jinghao Huang, Karimbergen Kudaybergenov, and Rui Liu,
\emph{The ball-covering property of von Neumann algebras and noncommutative
symmetric spaces}, arXiv:2311.14257.  It proves that a von Neumann algebra
has BCP/UBCP if and only if it is atomic and separably represented
(Theorem 3.8), and that a unital commutative \(C^*\)-algebra \(C(K)\) has
BCP/UBCP if and only if \(K=P(C(K))\) has a countable \(\pi\)-basis
(Corollary 3.9).  It then asks:

\begin{center}
\includegraphics[width=0.92\textwidth]{figures/open_problem_crop.png}
\end{center}

\section*{Transition Functionals}

For a \(C^*\)-algebra \(A\), define
\[
        T(A)=\ext(B_{A^*}),
\]
the extreme points of the dual unit ball, with the relative weak-star
topology inherited from \(A^*\).  These are the norm-one pure functionals.
In an irreducible representation they are the noncommutative analogue of
rank-one coefficient functionals
\[
        a\longmapsto \langle \pi(a)\xi,\eta\rangle .
\]
The point of this packet is to check that \(T(A)\), rather than only the
positive pure state space, matches the two boundary cases already solved in
the source paper.

\section*{Statement}

\begin{theorem}[commutative calibration]
\label{thm:comm}
Let \(K\) be a compact Hausdorff space.  Then \(T(C(K))\) has a countable
\(\pi\)-base if and only if \(K\) has a countable \(\pi\)-base.
Consequently, for a unital commutative \(C^*\)-algebra \(A=C(K)\), the
following are equivalent:
\[
        A\text{ has BCP},\qquad
        T(A)\text{ has a countable }\pi\text{-base},\qquad
        A\text{ has UBCP}.
\]
\end{theorem}

\begin{theorem}[von Neumann transition obstruction]
\label{thm:vn}
Let \(M\) be a von Neumann algebra.  If \(T(M)\) has a countable
\(\pi\)-base, then \(M\) is atomic and separably represented.  Consequently
\(M\) has BCP and UBCP by the source paper's Theorem 3.8.
\end{theorem}

\begin{corollary}
\label{cor:negative}
Every von Neumann algebra failing BCP in the source theorem has transition
functional space \(T(M)\) with no countable \(\pi\)-base.
\end{corollary}

\section*{Idea of the Proof}

For \(C(K)\), the Riesz representation theorem identifies the extreme points
of the dual unit ball with the point masses and phases:
\[
        T(C(K))\cong \T\times K,\qquad
        (\lambda,x)\mapsto \lambda\delta_x .
\]
Product with the metrizable circle preserves the existence of a countable
\(\pi\)-base.

For von Neumann algebras, the proof is a transition-functional version of the
pure-state obstruction.  In an atomless von Neumann algebra, the absolute
value of every pure functional is a singular pure state.  Given countably
many proposed transition-open sets, choose one pure functional from each.
A common projection kills all of their absolute values.  The transition-open
set of functionals that see that projection is nonempty, but none of the
chosen open sets is contained in it.

In the atomic but nonseparably represented case, uncountably many central
summands or uncountably many orthogonal rank-one projections give
uncountably many disjoint nonempty transition-open sets, ruling out a
countable \(\pi\)-base.

\section*{The Commutative Case}

\begin{proof}[Proof of Theorem~\ref{thm:comm}]
By the Riesz representation theorem, \(C(K)^*\) is the regular Borel measure
space \(M(K)\).  The extreme points of the unit ball of \(M(K)\) are exactly
\[
        \lambda\delta_x,\qquad \lambda\in\T,\ x\in K .
\]
The map \(\T\times K\to T(C(K))\), \((\lambda,x)\mapsto\lambda\delta_x\), is
a homeomorphism for the relative weak-star topology.

If \(K\) has a countable \(\pi\)-base \((U_n)\), and \((V_m)\) is a countable
base for \(\T\), then the sets
\[
        V_m\times U_n
\]
form a countable \(\pi\)-base for \(\T\times K\).

Conversely, suppose \(\T\times K\) has a countable \(\pi\)-base
\((W_n)\).  For every nonempty open \(U\subset K\), the set \(\T\times U\)
is a nonempty open subset of \(\T\times K\).  Choose \(n\) with
\[
        W_n\subset \T\times U .
\]
The projection of \(W_n\) onto \(K\) is a nonempty open subset of \(K\)
contained in \(U\).  Therefore the countable family of nonempty projections
of the \(W_n\)'s is a \(\pi\)-base for \(K\).

The final equivalence follows from the source paper's Corollary 3.9.
\end{proof}

\section*{A Common Projection Lemma}

\begin{lemma}[common projection for singular states]
\label{lem:common}
Let \(N\) be a nonzero von Neumann algebra and let
\((\varphi_n)\) be countably many singular states on \(N\).  Then there is a
nonzero projection \(p\in N\) such that
\[
        \varphi_n(p)=0\qquad(n\ge1).
\]
\end{lemma}

\begin{proof}
Fix a normal state \(\omega\) on \(N\).  For a singular positive functional
\(\varphi\), there are projections \(e_i\uparrow1\) with
\(\varphi(e_i)=0\).  Since \(\omega\) is normal, for every \(\varepsilon>0\)
there is a projection \(e\) with
\[
        \varphi(e)=0,\qquad \omega(1-e)<\varepsilon .
\]
Apply this to \(\varphi_n\) with \(\varepsilon_n=2^{-n-2}\), obtaining
projections \(e_n\) with
\[
        \varphi_n(e_n)=0,\qquad \omega(1-e_n)<2^{-n-2}.
\]
Set \(p=\bigwedge_n e_n\).  Then
\[
        \omega(1-p)
        \le
        \sum_n \omega(1-e_n)
        <\frac12,
\]
so \(p\ne0\).  Since \(p\le e_n\), positivity gives
\(\varphi_n(p)=0\) for every \(n\).
\end{proof}

\begin{lemma}
\label{lem:atomless-transition}
If \(N\) is a nonzero atomless von Neumann algebra, then \(T(N)\) has no
countable \(\pi\)-base.
\end{lemma}

\begin{proof}
Let \((U_n)\) be any countable family of nonempty open subsets of \(T(N)\).
Choose \(f_n\in U_n\).  For each \(n\), let \(|f_n|\) be the absolute value
of \(f_n\) in the polar decomposition of the functional.  Since \(f_n\) is an
extreme point of the dual unit ball, \(|f_n|\) is a pure state.  An atomless
von Neumann algebra has no normal pure states, so \(|f_n|\) is singular.

By Lemma~\ref{lem:common}, choose a nonzero projection \(p\in N\) such that
\[
        |f_n|(p)=0\qquad(n\ge1).
\]
The Cauchy--Schwarz inequality in the polar decomposition gives
\[
        |f_n(p)|^2\le |f_n|(p)=0,
\]
so \(f_n(p)=0\) for every \(n\).

Choose a pure state \(\psi_0\) on the nonzero unital \(C^*\)-algebra
\(pNp\).  Its canonical extension
\[
        \psi(x)=\psi_0(pxp)\qquad(x\in N)
\]
is a pure state of \(N\), and \(\psi(p)=1\).  Hence
\[
        W=\{f\in T(N): |f(p)|>1/2\}
\]
is a nonempty open subset of \(T(N)\).

But \(f_n\in U_n\setminus W\) for every \(n\).  Thus no \(U_n\) is contained
in \(W\), and the countable family \((U_n)\) cannot be a \(\pi\)-base.
\end{proof}

\section*{Atomic Nonseparable Von Neumann Algebras}

\begin{lemma}
\label{lem:disjoint}
If a topological space contains uncountably many pairwise disjoint nonempty
open subsets, then it has no countable \(\pi\)-base.
\end{lemma}

\begin{proof}
If \((V_n)\) were a countable \(\pi\)-base and \((O_i)_{i\in I}\) were
pairwise disjoint nonempty open sets, each \(O_i\) would contain some
\(V_{n(i)}\).  Distinct \(i\)'s require distinct \(n(i)\)'s, because the
\(O_i\)'s are disjoint and \(V_{n(i)}\ne\emptyset\).  Hence \(I\) would be
countable.
\end{proof}

\begin{lemma}
\label{lem:atomic-transition}
Let \(M\) be an atomic von Neumann algebra which is not separably represented.
Then \(T(M)\) has no countable \(\pi\)-base.
\end{lemma}

\begin{proof}
Write
\[
        M\cong \prod_{i\in I}\Bcal(H_i).
\]
If \(I\) is uncountable, let \(z_i\) be the central projection onto the
\(i\)-th summand.  The sets
\[
        O_i=\{f\in T(M): |f(z_i)|>1/2\}
\]
are nonempty and open.  They are pairwise disjoint: if \(i\ne j\), the
restriction of \(f\) to the abelian algebra generated by \(z_i,z_j\) has norm
at most one, so \(|f(z_i)|+|f(z_j)|\le1\).  Lemma~\ref{lem:disjoint} applies.

Now suppose \(I\) is countable but \(M\) is not separably represented.  Then
some summand is \(\Bcal(H)\) with \(H\) nonseparable.  Choose an uncountable
orthonormal family \((\xi_\alpha)_{\alpha\in\Gamma}\) in \(H\), and let
\(p_\alpha\) be the rank-one projection onto \(\C\xi_\alpha\).  The sets
\[
        W_\alpha
        =
        \{f\in T(\Bcal(H)): |f(p_\alpha)|>1/2\}
\]
are nonempty open subsets of \(T(\Bcal(H))\).  They are pairwise disjoint,
again because the restriction of a norm-one functional to the abelian
algebra generated by two orthogonal rank-one projections has norm at most
one.

The transition functional space \(T(\Bcal(H))\) sits as a nonempty open
subspace of \(T(M)\), determined by the corresponding central projection.
Thus \(T(M)\) also contains uncountably many pairwise disjoint nonempty open
sets.  Lemma~\ref{lem:disjoint} applies.
\end{proof}

\section*{Proof of the Von Neumann Calibration}

\begin{proof}[Proof of Theorem~\ref{thm:vn}]
We prove the contrapositive.  Suppose \(M\) is not atomic and separably
represented.

If \(M\) has a nonzero atomless central summand \(N=zM\), then \(T(N)\) is
identified with the open subspace
\[
        \{f\in T(M): |f(z)|>1/2\}.
\]
Indeed, an extreme norm-one functional is supported on exactly one side of
the central decomposition \(M=zM\oplus(1-z)M\).
By Lemma~\ref{lem:atomless-transition}, this open subspace has no countable
\(\pi\)-base, so \(T(M)\) has no countable \(\pi\)-base.

If \(M\) is atomic but not separably represented, then
Lemma~\ref{lem:atomic-transition} gives the same conclusion.

Therefore \(T(M)\) having a countable \(\pi\)-base forces \(M\) to be atomic
and separably represented.  The source paper's Theorem 3.8 then gives BCP
and UBCP.
\end{proof}

\begin{proof}[Proof of Corollary~\ref{cor:negative}]
The source theorem says that a von Neumann algebra fails BCP exactly when it
is not atomic and separably represented.  Apply Theorem~\ref{thm:vn}.
\end{proof}

\section*{Relation to the Full C*-Problem}

This packet is not a full solution of Question 3.  It says that the transition
functional space is a viable noncommutative replacement for the commutative
pure-state space at the two known boundaries:
\[
\begin{array}{c}
        C(K):\quad T(C(K))\cong \T\times K,\\[2mm]
        \text{von Neumann }M:\quad
        T(M)\text{ countable }\pi\text{-base}
        \Rightarrow M\text{ in the positive source class.}
\end{array}
\]

The missing arbitrary \(C^*\)-step is still a local bump theorem:
\[
        T(A)\text{ has enough countable local coefficient patches}
        \quad\Longrightarrow\quad
        A\text{ has UBCP}.
\]
The split-interval counterexample from the lane shows that such patches
cannot require separable principal right ideals.  The positive mechanism must
allow scalar or rank-one coefficient bumps over locally non-second-countable
sets.

\section*{Novelty and Search Bounds}

Local indexes were searched for arXiv:2311.14257, Corollary 3.9, Theorem
3.8, pure state space, pure transition functionals, extreme dual ball,
\(\pi\)-base, atomless von Neumann algebra, singular pure state, and the
previous lane-19 C*-BCP packets.  The source paper contains the von Neumann
BCP/UBCP classification and the commutative pure-state theorem, but no local
packet or source theorem recorded the transition-functional boundary
calibration in the form proved here.

\section*{Human Review Focus}

Reviewers should check:
\begin{enumerate}
    \item the homeomorphism \(T(C(K))\cong\T\times K\);
    \item the assertion that the absolute value of a pure functional is a
          pure state;
    \item the use of singularity for pure states on atomless von Neumann
          algebras;
    \item the Cauchy--Schwarz step \(|f(p)|^2\le |f|(p)\);
    \item the uncountable disjoint-open-set arguments for atomic
          nonseparable von Neumann algebras.
\end{enumerate}

\begin{thebibliography}{9}

\bibitem{HKL}
Jinghao Huang, Karimbergen Kudaybergenov, and Rui Liu,
\emph{The ball-covering property of von Neumann algebras and noncommutative
symmetric spaces},
arXiv:2311.14257, 2023.

\bibitem{Dixmier}
J. Dixmier,
\emph{\(C^*\)-Algebras},
North-Holland, 1977.

\bibitem{Sakai}
S. Sakai,
\emph{\(C^*\)-Algebras and \(W^*\)-Algebras},
Springer, 1971.

\bibitem{Blackadar}
Bruce Blackadar,
\emph{Operator Algebras: Theory of \(C^*\)-Algebras and Von Neumann
Algebras},
Encyclopaedia of Mathematical Sciences 122, Springer, 2006.

\bibitem{lane19state}
Run \texttt{fa\_banach\_001}, lane 19,
\emph{Pure-State Pi-Bases Obstruct the Negative Von Neumann Cases},
solution packet
\texttt{2311.14257\_von\_neumann\_pure\_state\_pi\_base\_obstruction}, 2026.

\bibitem{lane19split}
Run \texttt{fa\_banach\_001}, lane 19,
\emph{Pure-State Pi-Base Does Not Force a Separable Right-Support Net},
solution packet
\texttt{2311.14257\_pure\_pi\_base\_not\_support\_net\_split\_interval},
2026.

\end{thebibliography}

\end{document}
